\section{Generation}
\label{sec:generation}

\subsection{Model and runtime}

Answers are generated by Llama~3.2 with 3 billion parameters, served locally by
Ollama with a 4\,096-token context window, a temperature of 0 and JSON output
mode. At temperature 0 the same passages and question produce the same answer,
which makes behaviour reproducible and differences between versions of the
pipeline attributable to the pipeline. On the development machine about four
fifths of the model's layers fit in the 4\,GB of video memory and the rest run
from system memory, so generation speed is bounded by memory bandwidth rather
than by the GPU: an answer takes 7--9\,s once the model is loaded, and about
20\,s when it has to be loaded first. The interface warns in advance when the
model is not loaded.

The generator sits behind an abstract interface with a single method,
\texttt{generate(query, chunks, focus)}, so a different model or runtime can
be substituted by adding one class without touching retrieval or the API.

\paragraph{Trade-offs.} A local 3B model keeps every document on the machine,
costs nothing per question and works offline. The price is weaker synthesis
than a large hosted model and a small context window. Both push work into
retrieval: the passages must already be the right ones, in the right order,
and within budget, because the model cannot compensate for a poor selection.

\subsection{Prompt}

The prompt (\cref{lst:prompt}) states the model's role, gives it two tasks,
tells it what to do when the passages do not contain the answer, and fixes the
shape of its output. The \emph{answer} responds to the question directly; the
\emph{detail} states where in the document the answer comes from, citing
passages by number, and says plainly when the document only names something
without describing it.

\begin{lstlisting}[float, caption={The prompt, abridged
(\texttt{generation/basic\_backend.py}).}, label={lst:prompt}]
[The question names Article 101(1). The passages below are that
text; answer from them.]
You answer questions about a document using ONLY the numbered
context passages below.
Do two things:
1. ANSWER: answer the question directly and briefly.
2. DETAIL: say where in the document the answer comes from and what
   the surrounding text adds, citing passages by their number,
   e.g. [2]. If the document only names something without
   describing it, say so plainly instead of repeating the name.
Do not add anything the passages do not say.
If the passages do not contain the answer, say that they do not.
Respond with a JSON object EXACTLY in this shape:
{ "answer": "...", "detail": "..." }

Context:
[1] ...
[3] ...
[2] ...

Question: ...
\end{lstlisting}

\paragraph{Passage order.} Language models attend most reliably to the start
and end of a long context and least to its middle~\cite{liu2024lost}. The
passages arrive ranked, so they are interleaved before being placed in the
prompt: the first, third, fifth and so on from the front, then the rest in
reverse, which puts the strongest passages at both ends (for six passages the
order is 1, 3, 5, 6, 4, 2). Each passage keeps the number it had before
reordering, and that number is the one the interface shows, so ``[3]'' in the
detail is Source~3 in the citations pane.

\paragraph{Focus.} When exact lookup found the subdivision the question named,
the prompt opens with a line naming it: ``The question names Article 101(1).
The passages below are that text; answer from them.'' Without it, a small
model treats the passages as loose matches and may answer from a neighbouring
passage that reads as relevant.

\subsection{Deterministic replies}

Three situations are answered without relying on the model's judgement.

\paragraph{A named part that does not exist.} A question may name a
subdivision the provision does not have, such as paragraph 2 of an article
that is a single paragraph. Given the rest of the article, a small model
answers anyway, presenting a neighbouring sentence as the requested paragraph.
The model is therefore not called. The reply is composed from the stored
labels, and the whole provision is shown in the citations pane:

\begin{quote}\small
Article 85 has no paragraph 2: its text is not divided into paragraphs.\\
Article 7 has no paragraph 9. Its paragraphs are 1, 2, 3.
\end{quote}

\paragraph{A bare reference.} A query such as ``paragraph 1 article 101''
names a location but asks nothing, and with nothing to answer the model tends
to return an empty answer. A query is treated as a bare reference when nothing
but filler words (``show'', ``the text of'') remains once every recognised
reference is removed from it. The question given to the model is then replaced
by a task that makes the chat reply worth having next to the passages
themselves: ``Explain in plain language what Article 101(1) says.''

\paragraph{A non-answer.} If the model's answer is empty or a placeholder
(``None'', ``null'', ``N/A'', ``unknown''), it is replaced by a notice that
the passages are in the citations pane, rather than displayed as if it were
an answer.

\subsection{Response}

Every answer is returned as one JSON object: the answer and its detail; the
passages; for each passage, its document, pages and structural labels; the
cross-reference neighbourhood of the passages (\cref{sec:graph}); the lookup
note; retrieval and generation times; and the prompt and completion token
counts reported by Ollama together with the size of the context window.

\paragraph{No relevance percentage.} Earlier versions showed a relevance
percentage for each passage, derived from its cosine similarity to the
question. It was removed because it measured how closely the wording matched,
not whether the answer was right, and because passages found by exact lookup
always showed 100\,\%. The lookup note replaced it: the interface states
whether the passages are an exact match on the named provision or the closest
passages by keyword and meaning, which is what the reader needs in order to
judge them.
